\documentclass[11pt,a4paper]{article}

% ========================
% Packages
% ========================
\usepackage{geometry}
\usepackage{graphicx}
\usepackage{float}
\usepackage{hyperref}
\usepackage{fancyhdr}
\usepackage{array}
\usepackage{booktabs}
\usepackage{longtable}
\usepackage{enumitem}

\geometry{margin=1in}
\setlength{\parskip}{5pt}
\setlength{\parindent}{0pt}
\setlist{itemsep=2pt,topsep=3pt}

% Header/Footer
\pagestyle{fancy}
\fancyhf{}
\fancyhead[L]{End-Sem Report}
\fancyhead[R]{April, 2026}
\fancyfoot[C]{\thepage}

% ========================
% Title Page
% ========================
\begin{document}

\begin{titlepage}
\centering

{\Huge \textbf{\underline{End-Sem Report}}}\\[6em]

{\Huge \textbf{Atomic Fee Enforcement}}\\[2em]
{\Large \textbf{in Decentralized Perpetual Trading}}\\
\vspace{14px}
{\Large \textbf{A Transaction-Layer Approach Using Elliptic Curve Cryptography}}\\[5em]

\textbf{Course:} BLOCKCHAIN TECHNOLOGY (BITS F452) \\[0.5em]
\textbf{Instructor In-Charge:} Prof. Sanjay K. Sahay\\[5em]

\begin{center}
\setlength{\extrarowheight}{3pt}
{
\begin{tabular}{>{\centering\arraybackslash}p{0.55\textwidth} >{\centering\arraybackslash}p{0.35\textwidth}}
\textbf{Name of the Student} & \textbf{ID.No.} \\
Viyom Gupta & 2023A7PS0413G \\
\end{tabular}
}
\end{center}

\vspace{5em}

\includegraphics[width=0.3\textwidth]{bits.png}\\[2.5em]

\textbf{BIRLA INSTITUTE OF TECHNOLOGY \& SCIENCE, PILANI}\\[1em]
(April, 2026)

\end{titlepage}

% ========================
% Table of Contents
% ========================
\tableofcontents
\newpage

% ============================================================
% =============== ABSTRACT ====================================
% ============================================================
\section*{Abstract}
\addcontentsline{toc}{section}{Abstract}

Decentralized exchanges on Solana process billions of dollars in volume yet still rely on fragile fee-collection patterns that treat trade execution and platform-fee transfer as \textbf{two separate transactions} --- a structural source of revenue leakage, partial-execution bugs, and MEV-style race conditions. This project, \textbf{TRADEN-PROD}, eliminates the problem \textbf{without deploying any new on-chain program} by composing two existing primitives of the Solana runtime: single-signature coverage and atomic transaction execution. A custom interceptor in the application layer prepends a \texttt{SystemProgram.transfer} fee instruction into the user's trade transaction before it is signed, so the wallet's single Ed25519 signature over the SHA-256 hash of the message authorises both the fee and the trade as one indivisible unit.

The deliverable is a fully functional perpetual-futures trading platform (built on Drift Protocol) plus six purpose-built verification routes (\texttt{/blockchain}, \texttt{/verify}, \texttt{/security}, \texttt{/benchmarks}, \texttt{/explorer}, \texttt{/receipt/[signature]}) that allow any reviewer to \textit{click and verify} the cryptographic guarantees in real time. Six classes of attacks (fee evasion, fee tampering, recipient swap, replay, signature forgery, instruction reordering) were validated to fail; measured overhead is approximately 64 bytes per transaction, sub-millisecond extra signing latency, and at most 1\% of Solana's 200{,}000-CU compute budget.

% ---------- 1. Introduction ----------
\section{Introduction}

\subsection{Background \& Problem Statement}

Solana's high-throughput architecture (approximately 400\,ms slot time, \textasciitilde65{,}000 TPS theoretical capacity, \textasciitilde\$0.00025 per signature) has made it the dominant L1 for high-frequency on-chain trading. Drift Protocol, a perpetual futures DEX with 50+ markets, is one of the most active programs on Solana, and platforms built on top of it typically charge a 5--10 basis-point fee to be economically viable. That fee is almost universally collected as a separate transaction or as an off-chain payment, both of which assume the user (or a relayer) will \textit{cooperate} --- an unsafe assumption in a permissionless environment.

This separation produces three failure modes: \textbf{fee evasion} (a motivated user signs only the trade), \textbf{partial execution} (network congestion lands the trade while the fee expires), and \textbf{race conditions / MEV} (searchers reorder or front-run the fee). The deeper academic issue is that the design \textit{violates the principle of atomic execution} --- semantically related operations should succeed or fail together.

\subsection{Proposed Solution}

\textbf{Atomic Fee Enforcement} binds fee payment and trade execution into a single Solana transaction. The runtime guarantees every instruction commits or reverts together. Because the wallet produces a single Ed25519 signature over a SHA-256 hash of the entire serialized message, any modification --- removing the fee, lowering it, swapping the recipient, reordering instructions --- invalidates the signature and the validator rejects the transaction. Fee enforcement therefore becomes a \textit{cryptographic fact}, not a policy.

\subsection{Objectives Revisited}

\begin{center}
\begin{tabular}{c p{0.55\textwidth} l}
\toprule
\textbf{\#} & \textbf{Objective (from Mid-Sem)} & \textbf{Status} \\
\midrule
1 & Atomic Fee Bundling --- fee + trade in one all-or-nothing tx & Implemented \\
2 & Ed25519 Signing --- single signature authorises both & Implemented \\
3 & SHA-256 Integrity --- prevent post-signature tampering & Demonstrated \\
4 & Trade Primitive Support --- market, limit, stop-loss & All 5 order types \\
5 & V0 Transactions + Address Lookup Tables & Implemented \\
6 & Attack Resistance --- replay, forgery, bypass & 6 scenarios validated \\
\bottomrule
\end{tabular}
\end{center}

\subsection{Technology Stack}

Frontend: Next.js 15, React 19, TypeScript 5, Tailwind CSS 4, shadcn/ui. State: Zustand 5 (chain data), TanStack Query 5 (server data). Blockchain: Solana web3.js 1.98, Drift SDK 2.143, Solana Wallet Adapter (Phantom). In-browser cryptography: \texttt{tweetnacl} (Ed25519), Web Crypto API (SHA-256). Persistence: MongoDB + Mongoose 8. Charting: lightweight-charts. Oracles: Pyth Network and Switchboard.

% ---------- 2. Literature Survey ----------
\section{Literature Survey}

\subsection{Solana Transaction Model}

A Solana transaction is a \texttt{Message} containing one or more \texttt{Instruction}s and a \texttt{recent\_blockhash}. The runtime executes instructions in declared order and commits state changes only if every instruction returns \texttt{Ok}; otherwise all prior state changes revert. Versioned (V0) Transactions add support for Address Lookup Tables (ALTs), compressing 32-byte account references into single-byte indices and enabling far more complex instruction sets per transaction. The serialized message is signed with Ed25519 by every required signer.

\subsection{Cryptographic Foundations}

\textbf{Ed25519} (RFC 8032) is a deterministic Edwards-curve signature scheme over Curve25519 with 128-bit security, 32-byte keys, and 64-byte signatures. The single signature commits to every byte of the serialized message, so any tampering invalidates it; the scheme is also batch-verifiable, letting validators check thousands of signatures in parallel. \textbf{SHA-256} (FIPS 180-4) sits in the signing pipeline; its avalanche property (a one-bit input change flips $\sim$50\% of output bits) plus collision resistance ($\sim 2^{128}$ ops to forge) make tamper detection probabilistically certain.

\subsection{Drift Protocol \& Related Work}

Drift Protocol is Solana's largest decentralised perpetual exchange. Its \texttt{AuthorityDrift} SDK exposes \texttt{openPerpOrder} and \texttt{sendTransaction}, the natural injection points for fee prepending --- the entire enforcement engine is implemented as overrides on a freshly constructed \texttt{AuthorityDrift} instance.

\begin{center}
\begin{tabular}{l l l c}
\toprule
\textbf{Platform} & \textbf{Fee Mechanism} & \textbf{Atomicity Layer} & \textbf{New Program?} \\
\midrule
Jupiter & Fees inside swap program & Program-level & Yes \\
Raydium & LP fees in AMM math & Program-level & Yes \\
Drift Native & Builder Codes / RevenueShare & Protocol-level & Already deployed \\
\textbf{This work} & \textbf{Instruction prepending} & \textbf{Transaction-level} & \textbf{No} \\
\bottomrule
\end{tabular}
\end{center}

The transaction-layer approach is \textbf{portable} (works around any Solana DEX SDK) and incurs \textbf{zero on-chain audit cost} because no new program is deployed.

% ---------- 3. System Architecture ----------
\section{System Architecture}

The project organises into five layers: (1) presentation --- Next.js + React + Tailwind; (2) state \& sync --- Zustand stores and AuthorityDrift subscriptions for chain data, TanStack Query for server data; (3) \textbf{atomic fee enforcement} --- the project's contribution, implemented in \texttt{lib/DriftClientWrapper.ts} and \texttt{lib/tradingFee.ts}; (4) protocol \& wallet --- Drift SDK, Solana web3.js, Wallet Adapter, \texttt{tweetnacl}; (5) persistence \& APIs --- Next.js route handlers, MongoDB, Pyth/Switchboard oracles. Layer 3 is the only place that touches \texttt{sendTransaction}, so introducing a new market or screen does not require thinking about fees.

The application exposes thirteen routes. Trading routes: \texttt{/perps}, \texttt{/spot}, \texttt{/signals}, \texttt{/user}, \texttt{/data}. Wallet-gated dashboards: \texttt{/admin}, \texttt{/creator}. Education \& verification routes (added for end-sem): \texttt{/blockchain}, \texttt{/verify}, \texttt{/security}, \texttt{/benchmarks}, \texttt{/explorer}, \texttt{/receipt/[signature]}.

\textbf{Trade flow.} A click on \textbf{Place Order} invokes \texttt{drift.openPerpOrder(\ldots)}; the interceptor stashes the order parameters and delegates. The Drift SDK builds the trade transaction internally and calls \texttt{driftClient.sendTransaction(tx)}, which is also intercepted. The interceptor computes the fee in lamports (5\,bps; converted via SOL oracle if quote-asset), builds a \texttt{SystemProgram.transfer} instruction, decompiles the V0 message with all referenced ALTs resolved, prepends the fee instruction, recompiles to V0, and wraps it in a fresh \texttt{VersionedTransaction}. The wallet shows two instructions, the user approves once, and Phantom produces a single Ed25519 signature over \texttt{SHA-256(message)}. The validator verifies the signature, checks the recent-blockhash window, and executes both instructions atomically. On confirmation, \texttt{recordFeeToDatabase} writes a non-blocking MongoDB row and the UI links to \texttt{/receipt/[signature]}.

% ---------- 4. Implementation ----------
\section{Implementation}

\subsection{Atomic Fee Enforcement Engine (Core Innovation)}

Implemented in \texttt{ui/src/lib/DriftClientWrapper.ts}. The function \texttt{installTradingFeeInterceptor(drift)} is invoked once during \texttt{useSetupDrift} initialisation and applies two layered overrides on the active \texttt{AuthorityDrift} instance:

\begin{itemize}
    \item \textbf{\texttt{drift.openPerpOrder} override} captures \texttt{pendingPerpOrderFee = \{orderSize, assetType, marketIndex\}} into a closure, then delegates --- the context required to compute the fee correctly when \texttt{sendTransaction} is called moments later by SDK internals.
    \item \textbf{\texttt{driftClient.sendTransaction} override} computes the fee, builds a \texttt{SystemProgram.transfer} instruction, prepends it to the original transaction, and submits the modified transaction. The override is \textit{idempotent} (calling \texttt{installTradingFeeInterceptor} twice is a no-op) and \textit{fail-safe} (if fee construction throws, the trade is \textit{not} sent --- a fee can never be skipped silently).
\end{itemize}

The override branches on transaction type. For \textbf{V0 \texttt{VersionedTransaction}}, the original \texttt{MessageV0} is decompiled with all referenced ALTs resolved over RPC, the fee instruction is prepended, the message is recompiled to V0, and a fresh \texttt{VersionedTransaction} is constructed. Two non-obvious requirements were learned during implementation: ALTs must be resolved \textit{before} decompiling, and the same ALT array must be passed to \texttt{compileToV0Message} --- otherwise the validator rejects with \texttt{AccountNotFound}. For legacy \texttt{Transaction}, \texttt{transaction.instructions.unshift(feeInstruction)} suffices.

After successful submission, \texttt{recordFeeToDatabase} writes a \texttt{Fee} row to MongoDB containing \texttt{userId}, \texttt{experienceId}, \texttt{feeAmount}, \texttt{feeInLamports}, \texttt{orderSize}, \texttt{assetType}, and \texttt{txSignature}. This call is non-blocking and best-effort --- the canonical record of every fee remains the on-chain \texttt{SystemProgram.transfer}.

\subsection{Fee Computation \& Oracle-based USDC \texorpdfstring{$\rightarrow$}{->} SOL Conversion}

Implemented in \texttt{ui/src/lib/tradingFee.ts}. The platform fee is 5\,bps (0.05\%) of order size. For base-asset orders (e.g., size in SOL), \texttt{feeInLamports = orderSize $\times$ 5 / 10\,000} via \texttt{BN} arithmetic. For quote-asset orders (e.g., USDC), the USDC fee is converted to lamports using the current SOL price from Drift's oracle cache (Pyth + Switchboard): \texttt{feeInLamports = (feeInUsdc $\times$ 1e9) / solOraclePrice}. If the oracle is stale or unavailable, the trade is rejected --- never silently passed without a fee. The recipient \texttt{BUILDER\_AUTHORITY} is read from \texttt{NEXT\_PUBLIC\_BUILDER\_AUTHORITY}, allowing different deployments to route fees to different wallets without code changes.

\subsection{Trading Frontend \& Persistence}

The \texttt{/perps} page integrates a candlestick chart (lightweight-charts) with WebSocket price streaming, a live orderbook, positions and open-orders tables, and a trade form supporting all five Drift order types (market, limit, take-profit, stop-loss, oracle-limit) with a 1$\times$--20$\times$ leverage slider over 40+ markets. \texttt{/spot} provides deposit/withdraw/swap into Drift's spot pools. \texttt{/signals} executes creator-published signals through exactly the same path as manual trades, so the interceptor enforces the same atomic guarantee.

Four Mongoose schemas back the system: \texttt{FeeSchema} (off-chain mirror of the on-chain transfer, keyed by \texttt{txSignature}), \texttt{TradeSchema}, \texttt{SignalSchema}, and \texttt{FeeClaimSchema} (creator earning-claim workflow with PENDING $\rightarrow$ PROCESSING $\rightarrow$ COMPLETED transitions). API surface in \texttt{ui/src/app/api/}: \texttt{/api/fee}, \texttt{/api/trade}, \texttt{/api/signal}, \texttt{/api/fee-claim}, \texttt{/api/admin/stats}, \texttt{/api/admin/claims}.

\subsection{Education \& Verification Layer}

Six routes were added so every academic claim can be verified live:

\begin{itemize}
    \item \texttt{/blockchain} --- maps every BITS F452 syllabus topic to the exact file in the repo where it is implemented, with theory + code refs + live-demo links.
    \item \texttt{/verify} --- a 10-step in-browser wizard that generates a fresh Ed25519 keypair, builds a transaction, signs it, mutates one byte and shows the signature now fails. Includes a tamper matrix that flips every kind of edit.
    \item \texttt{/security} --- six adversarial simulations (replay, signature forgery, fee bypass, fee-amount manipulation, instruction reordering, MITM recipient swap), each emitting PASS/FAIL plus a byte-level diff.
    \item \texttt{/benchmarks} --- transaction-size delta, signing latency (100-iter avg/min/max), compute-unit estimate vs the 200{,}000-CU budget, end-to-end construction time.
    \item \texttt{/explorer} --- paste any signature; \texttt{getParsedTransaction} resolves it and the page highlights the fee instruction and the Drift instruction inside the same transaction.
    \item \texttt{/receipt/[signature]} --- per-trade card showing status, slot, blockhash, fee details, trade details, cryptographic proof block, and a Solscan deep-link.
\end{itemize}

All cryptographic demos run client-side using \texttt{Keypair.generate()} + \texttt{tweetnacl} + Web Crypto, so they require neither a connected wallet nor an RPC --- they work on any laptop during a presentation.

% ---------- 5. Concept Map ----------
\section{Cryptographic \& Blockchain Concepts --- Mapped to Code}

\begin{center}
\begin{tabular}{p{0.4\textwidth} p{0.5\textwidth}}
\toprule
\textbf{Concept (Course Syllabus)} & \textbf{Implementation in TRADEN-PROD} \\
\midrule
Asymmetric crypto (ECC) & Ed25519 over Curve25519, every Solana signature \\
Hash functions / SHA-256 & Solana signing pipeline + Web Crypto in \texttt{/verify} \\
Digital signatures & Phantom in production, \texttt{tweetnacl} in demos \\
Wallets, key derivation, addresses & Wallet Adapter, base58 pubkeys \\
Transactions \& atomicity & V0 \texttt{VersionedTransaction}; runtime atomicity exploited by the interceptor \\
V0 + Address Lookup Tables & \texttt{TransactionMessage.decompile} / \texttt{compileToV0Message} flow \\
Smart contracts & Drift Program + \texttt{SystemProgram} (no new program by us) \\
Decentralised oracles & Pyth + Switchboard via Drift oracle cache \\
Replay protection & 150-block recent-blockhash window, demoed in \texttt{/security} \\
Identity \& authentication & Pubkey-gated admin \& creator dashboards \\
Consensus (PoH + Tower BFT) & Solana finality surfaced via ``confirmed'' status in receipts \\
Scalability & Demonstrated on \texttt{/benchmarks} ($\le$ 1\% CU overhead) \\
\bottomrule
\end{tabular}
\end{center}

% ---------- 6. Security Analysis ----------
\section{Security Analysis}

\subsection{Threat Model}

The user is treated as adversarial and assumed to actively try to keep the trade while skipping the fee; the network may reorder transactions arbitrarily (MEV); the wallet correctly implements Ed25519 signing; and the Solana validator set behaves honestly per Tower BFT / PoH (the standard Solana trust assumption, not weakened here). No custom relayer, off-chain signing service, or platform-deployed Solana program is required.

\subsection{Attack-by-Attack Resistance}

\begin{center}
\begin{tabular}{c p{0.22\textwidth} p{0.5\textwidth} c}
\toprule
\textbf{\#} & \textbf{Attack} & \textbf{Defence} & \textbf{Status} \\
\midrule
1 & Fee Removal & Adversary deletes the fee instruction; new SHA-256 hash $\Rightarrow$ Ed25519 verification fails & Rejected \\
2 & Fee Amount Manipulation & Changing \texttt{lamports} field triggers SHA-256 avalanche $\Rightarrow$ signature invalid & Rejected \\
3 & Recipient Swap (MITM) & Account keys are part of the signed message; hash changes $\Rightarrow$ signature invalid & Rejected \\
4 & Replay Attack & \texttt{recent\_blockhash} part of signed message; expires after $\sim$150 slots & Rejected \\
5 & Signature Forgery & 128-bit Ed25519 security; $\sim 2^{128}$ ops infeasible & Infeasible \\
6 & Instruction Reordering & Order encoded in serialized bytes; hash changes $\Rightarrow$ signature invalid & Rejected \\
\bottomrule
\end{tabular}
\end{center}

\subsection{Why the System Is Cryptographically Indivisible}

Let $M = (\text{Fee\_instr}, \text{Trade\_instr}, \text{recent\_blockhash}, \text{account\_keys}, \ldots)$ be the serialized message and $\sigma = \text{Ed25519.sign}(\text{SHA-256}(M), sk)$. For the validator to accept the transaction it must verify $\sigma$ against the public key and the current serialised message. If an adversary delivers any $M' \ne M$, then with overwhelming probability $\text{SHA-256}(M') \ne \text{SHA-256}(M)$, hence $\text{Ed25519.verify}(\text{SHA-256}(M'), \sigma, pk) = \text{false}$. The Solana runtime then either commits every instruction in $M$ or none. Therefore Fee\_instr is committed \textbf{iff} Trade\_instr is committed --- fee evasion reduces to either forging an Ed25519 signature or finding a SHA-256 second pre-image, both computationally infeasible at the 128-bit security level.

% ---------- 7. Performance Analysis ----------
\section{Performance Analysis}

The \texttt{/benchmarks} page measures real overhead by constructing the same trade transaction in two variants (without fee, with fee) and comparing them. Adding one \texttt{SystemProgram.transfer} adds a 1-byte instruction-count delta, $\sim$32 bytes for the recipient account key, $\sim$9 bytes for the instruction data (4-byte transfer discriminator + 8-byte little-endian lamport amount), and small constant header overhead --- $\sim$64 bytes total out of Solana's 1{,}232-byte transaction limit. Ed25519 signing being a constant-time elliptic-curve operation, the additional in-browser signing latency over 100 iterations via \texttt{tweetnacl} is sub-millisecond (typically 0.2--0.5\,ms on commodity hardware). A \texttt{SystemProgram.transfer} consumes $\sim$150 compute units against Solana's default 200{,}000-CU budget, i.e. 0.075\%. The interceptor's decompile-prepend-recompile pipeline (including ALT resolution from RPC) adds 5--15\,ms on a normal devnet RPC, dominated by the network round-trip for \texttt{getAddressLookupTable}; on warm RPCs the additional time is negligible.

\begin{center}
\begin{tabular}{l l}
\toprule
\textbf{Metric} & \textbf{Result} \\
\midrule
Transaction-size overhead & $\sim$64 bytes ($\sim$5\% of 1{,}232-byte limit) \\
Signing latency overhead & sub-millisecond (100-iter avg) \\
Compute-unit overhead & $\sim$150 / 200{,}000 CU $=$ 0.075\% \\
Construction-time overhead & 5--15\,ms (RPC-bound) \\
\midrule
\textbf{Verdict} & \textbf{Atomic fee enforcement is essentially free.} \\
\bottomrule
\end{tabular}
\end{center}

% ---------- 8. Testing & Validation ----------
\section{Testing \& Validation}

The system was validated end-to-end on Solana devnet with a real Phantom wallet, a real Drift devnet deployment, and a real MongoDB Atlas cluster. Phantom integration (connection, public-key extraction, balance display, Ed25519 signing) was confirmed; test wallets were funded via the standard Solana faucet. Real perpetual positions were opened and closed across multiple markets; the platform fee instruction was visible in \textit{every} confirmed transaction on Solscan, and the same fact is exposed live via \texttt{/explorer} which calls \texttt{getParsedTransaction} to confirm both instructions are present in the same transaction. Every confirmed trade produces a \texttt{Fee} row and a \texttt{Trade} row whose \texttt{txSignature} matches the on-chain transaction. The six adversarial scenarios in \texttt{/security} were each attempted; all six produced the expected \texttt{INVALID} / \texttt{REJECTED} outcome. The four benchmarks were each run for 100 iterations and aggregated.

% ---------- 9. Discussion ----------
\section{Discussion: Limitations \& Honest Caveats}

\begin{itemize}
    \item \textbf{Drift Swift orders} bypass \texttt{driftClient.sendTransaction} (off-chain matching), so the interceptor is not exercised on Swift orders. The UI exposes a Swift toggle so the trade-off is visible to the user.
    \item \textbf{Application-layer enforcement} is binding only while the user is using TRADEN-PROD; the user is free to choose a different platform. This is a property of any application-layer fee model.
    \item \textbf{Off-chain database failures are tolerated.} A missing \texttt{Fee} row does not fail the on-chain trade --- the on-chain \texttt{SystemProgram.transfer} is the canonical record; the database is a convenience mirror for dashboards.
    \item \textbf{Oracle dependency} for USDC $\rightarrow$ SOL conversion: if both Pyth and Switchboard are simultaneously unavailable for the SOL feed, USDC-denominated trades fail safely rather than silently using a stale price.
\end{itemize}

% ---------- 10. Future Scope ----------
\section{Future Scope}

\begin{enumerate}
    \item \textbf{Mainnet deployment} with a production builder-authority wallet and signed cryptographic receipts attached to each fee row.
    \item \textbf{Cross-chain port} of the interceptor pattern to other runtimes that guarantee single-signature coverage and atomic execution (Sui, Aptos, Sei).
    \item \textbf{Dynamic fee engine} --- DAO-governed, tier-aware (volume discounts, loyalty rebates), still enforced atomically by the same prepend-instruction technique.
    \item \textbf{AI-driven trading signals} with on-chain commitments so creators stake reputation against published predictions.
    \item \textbf{Mobile wallet \& responsive layout}, plus packaging as a deployable creator-economy app on Whop.
\end{enumerate}

% ---------- 11. Conclusion ----------
\section{Conclusion}

TRADEN-PROD demonstrates that a real-world DeFi monetisation problem --- fee evasion in decentralised perpetual trading --- can be solved without deploying a new smart contract, without trusting a relayer, and without asking the user to behave well. The solution is a careful composition of two existing Solana primitives: the single-Ed25519-signature property and atomic transaction execution, glued together by an in-browser interceptor that prepends a \texttt{SystemProgram.transfer} instruction before signing. The cryptographic guarantee is mathematical: tampering with any byte of the signed message invalidates the signature with overwhelming probability under the security assumptions of Ed25519 and SHA-256. The performance cost is negligible: $\sim$64 bytes per transaction, sub-millisecond signing latency, and at most 1\% of the compute budget. Every major BITS F452 syllabus block --- asymmetric cryptography, hash functions, digital signatures, wallets, transactions, atomicity, replay protection, smart contracts, oracles, consensus, identity --- is implemented, mapped, demoed, and benchmarked inside the live application.

\begin{center}
\textit{Monetisation without trust. Trading without compromise.}
\end{center}

% ---------- 12. References ----------
\section{References}

\begin{enumerate}
    \item Y. Yakovenko, ``Solana: A New Architecture for a High Performance Blockchain,'' Solana Labs Whitepaper, 2018.
    \item Solana Foundation, ``Transaction Processing,'' \textit{Solana Documentation}, 2024. Available: \url{https://docs.solana.com/developing/programming-model/transactions}
    \item Solana Foundation, ``Versioned Transactions and Address Lookup Tables,'' 2024. Available: \url{https://docs.solana.com/developing/lookup-tables}
    \item S. Josefsson and I. Liusvaara, ``Edwards-Curve Digital Signature Algorithm (EdDSA),'' RFC 8032, IETF, 2017.
    \item National Institute of Standards and Technology, ``Secure Hash Standard (SHS),'' FIPS 180-4, 2015.
    \item Drift Labs, ``Drift Protocol Documentation,'' 2024. Available: \url{https://docs.drift.trade/}
    \item D. Boneh and V. Shoup, \textit{A Graduate Course in Applied Cryptography}, 2020. Available: \url{https://toc.cryptobook.us/}
    \item Phantom, ``Developer Documentation,'' 2024. Available: \url{https://docs.phantom.app/}
    \item Pyth Network, ``Pyth Price Feeds Documentation,'' 2024. Available: \url{https://docs.pyth.network/}
    \item Vercel, ``Next.js App Router Documentation,'' 2024. Available: \url{https://nextjs.org/docs}
\end{enumerate}

\end{document}
